\documentclass[UTF8,11pt,a4paper,fontset=none]{ctexart}

\setCJKmainfont[Path=C:/Windows/Fonts/,BoldFont=simhei.ttf,ItalicFont=STKAITI.TTF]{simsun.ttc}
\setCJKsansfont[Path=C:/Windows/Fonts/,BoldFont=msyhbd.ttc]{msyh.ttc}
\setCJKmonofont[Path=C:/Windows/Fonts/]{msyh.ttc}

\usepackage[a4paper,left=2.35cm,right=2.35cm,top=2.25cm,bottom=2.3cm,
  headheight=22pt,headsep=0.55cm,footskip=0.9cm]{geometry}
\usepackage{amsmath,amssymb,mathtools,bm}
\usepackage{booktabs,tabularx,array,multirow,colortbl,makecell}
\usepackage{enumitem}
\usepackage{tikz}
\usetikzlibrary{arrows.meta,positioning,calc,shapes.geometric}
\usepackage[most]{tcolorbox}
\usepackage{xcolor}
\usepackage{hyperref}
\usepackage{cleveref}
\usepackage{caption}
\usepackage{subcaption}
\usepackage{multicol}
\usepackage{fancyhdr}
\usepackage{lastpage}
\usepackage{titlesec}
\usepackage{needspace}
\usepackage{setspace}
\usepackage[nopatch=footnote]{microtype}
\usepackage{xurl}
\usepackage{newtxmath}

\definecolor{ink}{HTML}{202124}
\definecolor{navy}{HTML}{17365D}
\definecolor{blue}{HTML}{255C8B}
\definecolor{teal}{HTML}{267A78}
\definecolor{orange}{HTML}{B85C24}
\definecolor{green}{HTML}{477A52}
\definecolor{red}{HTML}{9D3545}
\definecolor{gray}{HTML}{647184}
\definecolor{line}{HTML}{D9E0E8}
\definecolor{graybg}{HTML}{F7F7F7}
\definecolor{lightblue}{HTML}{F4F7FA}
\definecolor{lightgreen}{HTML}{F5F8F5}
\definecolor{lightorange}{HTML}{FAF7F2}
\definecolor{paperblue}{HTML}{FAFBFC}

\hypersetup{colorlinks=true,linkcolor=navy,citecolor=navy,urlcolor=blue,
  pdftitle={ProRM: 局部策略效用、Fisher 修正与单步 NGD},
  pdfauthor={ProRM Research Report}}

\pagestyle{fancy}
\fancyhf{}
\fancyhead[L]{\footnotesize\color{gray}\textsc{Prospective Reward Modeling}}
\fancyhead[R]{\footnotesize\color{gray}局部策略效用与奖励等价类}
\fancyfoot[C]{\footnotesize\color{gray}\thepage}
\renewcommand{\headrulewidth}{0.3pt}
\renewcommand{\headrule}{\hbox to\headwidth{\color{line}\leaders\hrule height \headrulewidth\hfill}}
\renewcommand{\footrulewidth}{0pt}

\setlength{\parindent}{2em}
\setlength{\parskip}{0.12em}
\setstretch{1.16}
\setlength{\abovedisplayskip}{0.65em}
\setlength{\belowdisplayskip}{0.65em}
\setlength{\abovedisplayshortskip}{0.4em}
\setlength{\belowdisplayshortskip}{0.4em}
\setlist[itemize]{leftmargin=1.6em,itemsep=0.28em,topsep=0.4em,parsep=0.05em}
\setlist[enumerate]{leftmargin=1.9em,itemsep=0.34em,topsep=0.4em,parsep=0.05em}
\captionsetup{font={small,stretch=1.08},labelfont={bf},labelsep=quad}
\renewcommand{\arraystretch}{1.22}
\allowdisplaybreaks

\titleformat{\section}[hang]
  {\Needspace{5\baselineskip}\sffamily\Large\bfseries\color{navy}}
  {\thesection}{0.8em}{}
\titleformat{\subsection}[hang]
  {\Needspace{4\baselineskip}\sffamily\large\bfseries\color{ink}}
  {\thesubsection}{0.7em}{}
\titlespacing*{\section}{0pt}{2.1em}{0.9em}
\titlespacing*{\subsection}{0pt}{1.45em}{0.55em}

\newcommand{\R}{\mathbb{R}}
\newcommand{\E}{\mathbb{E}}
\newcommand{\KL}{D_{\mathrm{KL}}}
\newcommand{\Reg}{\operatorname{Reg}}
\newcommand{\Null}{\operatorname{Null}}
\newcommand{\row}{\operatorname{row}}
\newcommand{\diag}{\operatorname{diag}}
\newcommand{\sig}{\sigma}
\newcommand{\argmin}{\operatorname*{arg\,min}}
\newcommand{\argmax}{\operatorname*{arg\,max}}
\newcommand{\rstar}{r^{\star}}
\newcommand{\piref}{\pi_0}
\newcommand{\Fisher}{F_0}
\newcommand{\LocalRegret}{\widetilde{\Reg}}
\newcommand{\term}[1]{\textcolor{navy}{\sffamily\bfseries #1}}

\newtcolorbox{keybox}[1][]{enhanced,colback=lightblue,colframe=line,
  boxrule=0.5pt,borderline west={1.2pt}{0pt}{navy},arc=0mm,
  left=3mm,right=3mm,top=2.1mm,bottom=2.1mm,
  before skip=0.8em,after skip=0.8em,fonttitle=\sffamily\bfseries\color{navy},title=#1}
\newtcolorbox{claimbox}[1][]{enhanced,colback=lightgreen,colframe=line,
  boxrule=0.5pt,borderline west={1.2pt}{0pt}{green},arc=0mm,
  left=3mm,right=3mm,top=2.1mm,bottom=2.1mm,
  before skip=0.8em,after skip=0.8em,fonttitle=\sffamily\bfseries\color{ink},title=#1}
\newtcolorbox{warnbox}[1][]{enhanced,colback=lightorange,colframe=line,
  boxrule=0.5pt,borderline west={1.2pt}{0pt}{orange},arc=0mm,
  left=3mm,right=3mm,top=1.5mm,bottom=1.5mm,fontupper=\small,
  before skip=0.8em,after skip=0.8em,fonttitle=\sffamily\small\bfseries\color{ink},title=#1}
\newtcolorbox{definitionbox}[1][]{enhanced,colback=paperblue,colframe=line,
  boxrule=0.55pt,arc=1mm,left=3mm,right=3mm,top=2mm,bottom=2mm,
  before skip=0.7em,after skip=0.7em,fonttitle=\sffamily\bfseries\color{navy},title=#1}

\begin{document}
\color{ink}
\begin{center}
{\sffamily\fontsize{22}{28}\selectfont\bfseries\color{navy}
Prospective Reward Modeling}\par
\vspace{0.45em}
{\sffamily\Large\bfseries 面向下游策略效用的奖励学习}\par
\vspace{0.45em}
{\large 局部后悔、奖励等价类与 Fisher 修正的单步自然梯度}\par
\vspace{0.9em}
{\small\color{gray}理论与实验报告\qquad \today}
\end{center}
\vspace{0.8em}
{\color{line}\rule{\linewidth}{0.7pt}}

\begin{abstract}
标准 BT-RM 通过最大似然拟合历史偏好标签；但在受限奖励类和受限策略类下，偏好概率拟合得好并不必然导致下游策略效用最优。本文将奖励学习重新定义为 fixed reference policy、fixed policy tangent、fixed one-step optimizer 下的下游效用最大化问题。全局效用是双层优化，难以直接训练；在参考策略附近作 Taylor--Fisher 局部化后，得到闭式目标
\[
\widetilde{\Reg}(r)=\frac{1}{2\beta}\|A_0(r-r^*)\|_{F_0^\dagger}^2.
\]
该式刻画的是概率加权 reward error 在 policy-visible 空间的正交投影，而不是两个 Bernoulli 偏好分布的 KL。由此产生奖励等价类：若受限奖励类与真实奖励等价类相交，则无需恢复真实奖励本身；若无交集，则应选择局部近似下效用最高的可实现等价类。为落地到 pairwise preference 数据，ProRM 使用同一 pair 的重复条件独立标注构造无偏 margin 统计量，并通过 Fenchel 对偶得到可训练的 Fisher--GMM min--max 目标。有限样本实现用 train-only prompt cross-fit 选择跨 seed、跨 reward source 共享的 Fisher 正则逆，并在同一 $\beta$ 下比较 MLE-RM、ProRM 与 oracle 所诱导的单步自然梯度策略。在 $\beta\in\{0.1,0.2,0.3\}$ 上，三 seed 聚合均给出 $J_{\rm tab}>J_{\rm Pro}>J_{\rm oracle}>J_{\rm MLE}>J_{\pi_0}$，且 Pro 在每个 seed 上均优于 MLE 与 oracle；与此同时，ProRM 的 held-out NLL、奖励 MSE 与局部近似 regret 均差于 MLE-RM，揭示了清晰的 reward-model 泛化缺口。本文据此将奖励层诊断与下游 policy 结论严格分开。
\end{abstract}

\noindent\textbf{关键词：} 奖励建模；策略效用；奖励等价类；Fisher 几何；自然梯度

\subsection*{主要记号}
\small
\begin{tabularx}{\linewidth}{>{\bfseries\color{navy}}p{1.8cm} X >{\bfseries\color{navy}}p{1.8cm} X}
$r^*,\ r_\phi$ & 真实/操作性 oracle 奖励与受限奖励模型 & $\pi_0,\ \pi_\theta$ & 参考策略与更新后策略 \\
$A_0$ & reward 误差到 policy tangent 的局部映射 & $F_0$ & $\pi_0$ 处的 population Fisher 几何 \\
$\beta$ & 理论与实验 $J$ 中 forward-KL 的共同惩罚系数 & $\widehat F_\lambda$ & train-only 阻尼经验 Fisher \\
$\lambda_{\rm rel}$ & 经验 Fisher 逆的相对 damping & $\delta$ & 固定 LoRA-$B$ tangent 中的一步参数更新 \\
\end{tabularx}
\normalsize

% ================================================================
\section{局部策略效用与奖励等价类}

\subsection{问题设定：偏好拟合与下游效用的错位}

把 prompt--response 支持上的真实奖励写成 $r^*\in\R^m$，局部策略参数写成 $\theta\in\R^d$，奖励模型参数写成 $\phi\in\R^p$。在 LLM 后训练中，典型关系是
\[
\boxed{m\gg d\gg p.}
\]
$m\gg d$ 表示策略不可能利用 reward 空间所有方向；$d\gg p$ 表示受限奖励模型又无法表达所有 policy-visible 方向。因此，最小化 reward pointwise error 或 BT likelihood 都不是天然正确的目标。正确对象应是：
\[
\text{用 }r_\phi\text{ 训练出来的策略，在真实 }r^*\text{ 下的效用。}
\]

这与 decision-aware 的 ``predict, then optimize'' 原则一致：预测误差应由其诱导的决策损失评价 \cite{spo}。已有 RLHF 理论研究了 pairwise/$K$-wise reward MLE、pessimism 及两阶段 RLHF 与 DPO 的统计差异 \cite{zhu_principled,nika_rml_dpo,li_rl_dpo}；近期工作进一步刻画了 representation misspecification、条件等价与局部自然梯度结构 \cite{shi_gap,yang_equivalence,gopalan_misspecified}。本文保留显式 reward model，聚焦一个不同问题：受限奖励类中，应该学习哪一个对下游 policy 有效的奖励代表元。

给定参考策略 $\pi_0=\pi_{\theta_0}$ 和 KL 系数 $\beta>0$，并预先固定包含初始化、随机性与 tie-break 的下游优化器选择映射 $\mathcal A$，定义
\begin{align}
J(\theta;r)&=\E_{x\sim\rho,\,y\sim\pi_\theta(\cdot\mid x)}[r(x,y)]
-\beta\E_{x\sim\rho}\KL\bigl(\pi_\theta(\cdot\mid x)\,\|\,\pi_0(\cdot\mid x)\bigr),\\
\theta_r&=\mathcal A(r)\in\argmax_\theta J(\theta;r),\qquad U(r)=J(\theta_r;r^*),\\
\Reg(r)&=U(r^*)-U(r).
\end{align}
这是最自然的下游效用定义，但 $r\mapsto\theta_r\mapsto U(r)$ 是全局双层优化，不适合作为第一阶段监督训练目标。

\subsection{Taylor--Fisher 局部化与闭式后悔}

令
\[
d_0(x,y)=\rho(x)\pi_0(y\mid x),\qquad
s_0(x,y)=\nabla_\theta\log\pi_{\theta_0}(y\mid x).
\]
把所有 score 作为列组成 $A_{\theta_0}\in\R^{d\times m}$，并定义
\[
D_0=\diag(d_0),\qquad A_0=A_{\theta_0}D_0,
\qquad F_0=A_{\theta_0}D_0A_{\theta_0}^{\top}=\E_{d_0}[s_0s_0^\top].
\]
这里 $F_0$ 是参考策略处的 Fisher 几何 \cite{amari}。
在 $\theta_0$ 附近令 $\delta=\theta-\theta_0$。对 reward expectation 保留一阶、对 KL 保留二阶，则 quadratic surrogate 为
\[
J(\theta;r)\approx C(r)+\delta^\top A_0r-\frac{\beta}{2}\delta^\top F_0\delta.
\]
局部最优方向为
\[
\boxed{\delta_r=\frac1\beta F_0^\dagger A_0r.}
\]
代入真实奖励 $r^*$ 的局部效用并完成平方，得到
\begin{equation}
\boxed{
\widetilde{\Reg}(r)=\frac{1}{2\beta}
\bigl(A_0(r-r^*)\bigr)^\top F_0^\dagger
\bigl(A_0(r-r^*)\bigr).
}
\label{eq:local-regret}
\end{equation}
这就是 ProRM 的理想 population surrogate。它不是任意全局最优点上的恒等式；若固定优化器停留在同一光滑 local branch 且 $\delta=O(\beta^{-1})$，则原非线性 regret 的 leading term 为 \cref{eq:local-regret}，余项为 $O(\beta^{-2})$。

\subsection{奖励等价类：从真实奖励到等价类}

定义局部奖励等价关系
\[
r_1\sim r_2\quad\Longleftrightarrow\quad A_0r_1=A_0r_2.
\]
等价类为
\[
[r]=r+\Null(A_0).
\]
同一等价类中的所有奖励诱导相同的局部 surrogate 策略方向和相同的 leading-order 局部效用。于是 realizability 应分成三层：
\begin{enumerate}
\item \textbf{完整可实现：} $r^*\in\mathcal R$。BT-RM 与 ProRM 在 population 下可一致。
\item \textbf{等价类可实现：} $r^*\notin\mathcal R$，但 $\mathcal R\cap[r^*]\neq\varnothing$。此时无需恢复 $r^*$，只需命中真实等价类。
\item \textbf{等价类无交集：} $\mathcal R\cap[r^*]=\varnothing$。此时应选择
\[
r_{\mathrm{Pro}}\in\argmin_{r\in\mathcal R}\|A_0(r-r^*)\|_{F_0^\dagger}^2,
\]
即局部 surrogate 效用最高的可实现等价类。
\end{enumerate}

\begin{claimbox}[命题：局部奖励等价类]
两阶段局部 RLHF 的核心不是 ``学习真实奖励向量''，而是 ``学习真实奖励在 policy quotient space 中的像''。这将可实现条件从 $r^*\in\mathcal R$ 放松为 $\mathcal R\cap[r^*]\neq\varnothing$。
\end{claimbox}

\subsection{几何解释：正交投影与 Bernoulli--KL 投影}

令 $B_0=A_{\theta_0}D_0^{1/2}$，$h=r-r^*$。由于 $F_0=B_0B_0^\top$，有
\begin{align}
\|A_0h\|_{F_0^\dagger}^{2}
&=h^\top D_0^{1/2}B_0^\top(B_0B_0^\top)^\dagger B_0D_0^{1/2}h\\
&=\left\|P_{\row(B_0)}D_0^{1/2}h\right\|_2^2.
\end{align}
因此
\begin{equation}
\boxed{
\widetilde{\Reg}(r)=\frac{1}{2\beta}
\left\|P_{\row(A_{\theta_0}D_0^{1/2})}D_0^{1/2}(r-r^*)\right\|_2^2.
}
\label{eq:projection}
\end{equation}
这说明 ProRM 最小化的是：概率加权 reward error 在 policy-visible row space 上的正交投影长度。

在 Bradley--Terry 模型 \cite{bt} 下，BT-RM 的 population excess NLL（即减去与 $r$ 无关的 truth entropy 后的部分）为
\begin{equation}
\mathcal D_{\mathrm{BT}}(r)=\E_{e\sim Q_0}
 d_{\mathrm{KL}}\!\bigl(
\operatorname{Bern}(\sig(\Delta r^*(e)))\,\|\,
\operatorname{Bern}(\sig(\Delta r(e)))
\bigr),
\label{eq:bt-kl}
\end{equation}
其中 $Q_0=\rho\pi_0\pi_0$。它会拟合所有 preference-visible 方向，包括 policy-invisible 的 $\Null(A_0)$ 成分。因此，在受限奖励类中，\cref{eq:projection} 与 \cref{eq:bt-kl} 一般给出不同投影。

% ================================================================
\section{可解析例证：奖励等价类与投影错位}

\subsection{模型设定}

有三个等概率 context $x_1,x_2,x_3$，每个 context 有两个 response $\{+,-\}$，reference policy 均匀。策略只有二维局部参数：
\[
\pi_\theta(+\mid x_j)=\sig(2z_j^\top\theta),
\quad z_1=e_1,\\ z_2=e_2,\\ z_3=0,
\quad \theta_0=0.
\]
因此 $x_1,x_2$ 是 policy-visible，$x_3$ 是 preference-visible 但 policy-invisible：它能出现在偏好数据里，却无法被当前 policy tangent 改变。

用 margin $d_j=r(x_j,+)-r(x_j,-)$ 表示奖励，取中心化代表 $r_d(x_j,\pm)=\pm d_j/2$。则
\begin{equation}
A_0r_d=\frac16(d_1,d_2)^\top,
\qquad
F_0=\frac13I_2.
\label{eq:simple-AF}
\end{equation}
所以奖励等价类只由
\[
\boxed{g(d)=(d_1,d_2)}
\]
决定；$d_3$ 完全在 $\Null(A_0)$ 中。

令受限奖励类为一维共享特征
\[
\boxed{\mathcal R=\{d(w)=(w,w,w):w\in\R\}.}
\]
直观解释：奖励模型只有一个全局强度参数 $w$，必须对三个 context 给出同一个 margin；但策略实际只用得上前两个 margin。

\subsection{BT-RM 与 ProRM 的闭式解}

令
\[
\ell(z,t)=\log(1+e^t)-\sig(z)t.
\]
完整自然 $Q_0$ 下，BT population NLL 为
\begin{equation}
L_{\mathrm{BT}}(w)=\frac12\log2+\frac16\sum_{j=1}^{3}\ell(d_j^*,w),
\label{eq:example-nll}
\end{equation}
所以
\begin{equation}
\boxed{
w_{\mathrm{BT}}=\
\operatorname{logit}\left(\frac{\sig(d_1^*)+\sig(d_2^*)+\sig(d_3^*)}{3}\right).
}
\label{eq:example-wbt}
\end{equation}
BT-RM 会被三个 context 的 preference probability 共同决定，包括 policy-invisible 的 $d_3^*$。

ProRM 只看 $g(d)=(d_1,d_2)$。由 \cref{eq:simple-AF}，
\begin{equation}
\boxed{
\widetilde{\Reg}(w)=\frac{(w-d_1^*)^2+(w-d_2^*)^2}{24\beta},
\qquad
w_{\mathrm{Pro}}=\frac{d_1^*+d_2^*}{2}.
}
\label{eq:example-wpro}
\end{equation}
ProRM 完全忽略 $d_3^*$，因为 $x_3$ 对当前局部策略不可控。

\subsection{三类可实现性情形}

固定 $\beta=8$。三个 case 只改变真实 margin $d^*$，其余设置完全相同。

\begin{table}[ht]
\centering
\small
\setlength{\tabcolsep}{4pt}
\begin{tabularx}{\textwidth}{>{\bfseries}c X c c c}
\toprule
\rowcolor{lightblue}
Case & 几何含义 & $d^*$ & $w_{\mathrm{BT}}$ & $w_{\mathrm{Pro}}$\\
\midrule
I & 完整真实奖励可实现：$d^*=d(2)\in\mathcal R$ & $(2,2,2)$ & $2.000$ & $2.000$\\
II & 等价类可实现：$d^*\notin\mathcal R$，但 $d(2)=(2,2,2)\in\mathcal R\cap[d^*]$ & $(2,2,6)$ & $2.438$ & $2.000$\\
III & 等价类无交集：真实 visible 点 $(6,0)$ 不在直线 $(w,w)$ 上 & $(6,0,-6)$ & $0.000$ & $3.000$\\
\bottomrule
\end{tabularx}
\caption{同一最小模型下的三个 case。Case II 中，policy-invisible 的强偏好 $d_3^*=6$ 会拉动 BT-RM，却不应改变策略；Case III 中，ProRM 选择 visible quotient space 中的正交投影。}
\label{tab:cases}
\end{table}

\begin{figure}[ht]
\centering
\begin{subfigure}[t]{0.46\textwidth}
\centering
\begin{tikzpicture}[scale=0.70,>=Latex]
  \draw[->] (-0.25,0)--(5.4,0) node[right] {$g_1=d_1$};
  \draw[->] (0,-0.25)--(0,5.4) node[above] {$g_2=d_2$};
  \draw[thick,navy] (0,0)--(5.15,5.15) node[above right] {$\mathcal G:(w,w)$};
  \foreach \x in {1,2,3,4,5}{\draw (\x,0.06)--(\x,-0.06) node[below,font=\scriptsize]{\x};\draw (0.06,\x)--(-0.06,\x) node[left,font=\scriptsize]{\x};}
  \filldraw[green] (2,2) circle (3pt) node[below right,font=\scriptsize] {ProRM: $w=2$};
  \filldraw[orange] (2.438,2.438) rectangle +(0.14,0.14) node[above right,font=\scriptsize] {BT: $w=2.438$};
  \node[font=\small\bfseries] at (2.55,5.1) {Case II：等价类有交集};
  \node[align=center,font=\scriptsize] at (3.2,1.0) {$d_3^*=6$ 只在 null 方向\\却会拉动 BT};
\end{tikzpicture}
\end{subfigure}\hfill
\begin{subfigure}[t]{0.46\textwidth}
\centering
\begin{tikzpicture}[scale=0.70,>=Latex]
  \draw[->] (-0.25,0)--(6.7,0) node[right] {$g_1=d_1$};
  \draw[->] (0,-0.25)--(0,6.0) node[above] {$g_2=d_2$};
  \draw[thick,navy] (0,0)--(5.6,5.6) node[above right] {$\mathcal G:(w,w)$};
  \foreach \x in {1,2,3,4,5,6}{\draw (\x,0.06)--(\x,-0.06) node[below,font=\scriptsize]{\x};\draw (0.06,\x)--(-0.06,\x) node[left,font=\scriptsize]{\x};}
  \filldraw[red] (6,0) circle (3pt) node[above right,font=\scriptsize] {真实 visible $(6,0)$};
  \filldraw[green] (3,3) circle (3pt) node[above left,font=\scriptsize] {ProRM: $w=3$};
  \filldraw[orange] (0,0) rectangle +(0.14,0.14) node[below right,font=\scriptsize] {BT: $w=0$};
  \draw[dashed,green,->] (5.65,0.25)--(3.25,2.75);
  \node[font=\small\bfseries] at (3.2,5.85) {Case III：等价类无交集};
\end{tikzpicture}
\end{subfigure}
\caption{最直观的 quotient-space 图。横纵轴是 policy 能改变的两个 context；$d_3$ 不画出，因为它属于 $\Null(A_0)$。BT-RM 会被 $d_3$ 影响，ProRM 不会。}
\label{fig:simple-example}
\end{figure}

\begin{table}[ht]
\centering
\small
\setlength{\tabcolsep}{3.8pt}
\begin{tabular}{ccrrrrrr}
\toprule
\rowcolor{lightblue}
\multirow{2}{*}{Case} & \multirow{2}{*}{方法} & \multicolumn{2}{c}{自然 $Q_0$ NLL $\downarrow$} & \multicolumn{2}{c}{局部 regret $\downarrow$} & \multicolumn{2}{c}{真实 regret $\downarrow$}\\
\cmidrule(lr){3-4}\cmidrule(lr){5-6}\cmidrule(lr){7-8}
&& 数值 & 优者 & 数值 & 优者 & 数值 & 优者\\
\midrule
I & BT-RM & 0.5292 & \checkmark & 0 & \checkmark & 0 & \checkmark\\
I & ProRM & 0.5292 & \checkmark & 0 & \checkmark & 0 & \checkmark\\
\midrule
II & BT-RM & \textbf{0.4863} & \checkmark & 0.00200 & & 0.00196 & \\
II & ProRM & 0.4903 & & \textbf{0} & \checkmark & \textbf{0} & \checkmark\\
\midrule
III & BT-RM & \textbf{0.6931} & \checkmark & 0.1875 & & 0.1833 & \\
III & ProRM & 1.1209 & & \textbf{0.0938} & \checkmark & \textbf{0.0901} & \checkmark\\
\bottomrule
\end{tabular}
\caption{三 case 的 population 结果。Case II：BT-RM 有更低 NLL，但 ProRM 达到零 regret。Case III：ProRM 的 NLL 明显更差，但将局部 regret 降低 50\%，真实 regret 降低约 51\%。}
\label{tab:results}
\end{table}

真实 regret 在原始非线性 KL 正则目标下计算：
\[
\Reg_{\mathrm{true}}(w)=\frac\beta3\left[
\KL\bigl(\operatorname{Bern}(\sig(w/\beta))\,\|\,\operatorname{Bern}(\sig(d_1^*/\beta))\bigr)+
\KL\bigl(\operatorname{Bern}(\sig(w/\beta))\,\|\,\operatorname{Bern}(\sig(d_2^*/\beta))\bigr)
\right].
\]
$x_3$ 不出现，因为当前 policy tangent 无法改变它。这一例证得到如下结论：

\begin{claimbox}[例证结论]
policy-invisible 的偏好信号可以显著改善 BT likelihood，却不应该改变下游策略；当 reward class 受限时，BT-RM 可能为了拟合这些无效偏好而选择错误等价类，ProRM 则直接在 policy-visible quotient space 中做效用投影。
\end{claimbox}

% ================================================================
\section{ProRM 估计与 Fisher 修正的单步 NGD 主实验}

\subsection{偏好数据、无偏边际统计量与 ProRM 估计}

ProRM 使用两类数据对象。
\begin{itemize}
\item \textbf{节点数据：} $(x,y)\sim d_0$，用于估计 Fisher。工程中保存分数行 $s_0(x,y)$，形成矩阵 $S$，并以
\[
\widehat F_0u=\frac1{n_F}S^\top(Su)
\]
计算 Fisher--向量乘积。
\item \textbf{边数据：} $e=(x,y,y')\sim Q_0$，令
\[
z_0(e)=s_0(x,y)-s_0(x,y'),\qquad t_\phi(e)=\Delta r_\phi(e).
\]
在同一 $e$ 上采集条件独立的重复标注。若 $a_j\mid e\sim\operatorname{Bernoulli}(\sig(\Delta r^*(e)))$，与标注独立的 $N$ 满足 $\Pr(N=n)=(1-\gamma)\gamma^{n-1}$（$n\geq1$），且 $S_N=\sum_{j=1}^Na_j$，定义
\[
H(e)=\sum_{k=1}^{N}\frac{\binom{S_N}{k}-\binom{N-S_N}{k}}{k\gamma^{k-1}\binom Nk}.
\]
则 $\E[H(e)\mid e]=\Delta r^*(e)$。实验锁定 $\gamma=0.9$，并通过 oracle 变换保证 $p^*\in[0.25,0.75]$，因而该估计量具有有限二阶矩；本文不进一步声称轻尾性或有限四阶矩。
\end{itemize}

更一般地，比较对的选择会改变下游信息几何与样本效率 \cite{han_pairs}。本实验固定候选 $0$--$1$ 的规范边，目的是隔离 BT-MLE 与 ProRM 的目标差异，而非声称该比较对设计最优。

由配对恒等式可得
\[
A_0r=\frac12\E_{Q_0}[z_0(e)\Delta r(e)],\qquad
F_0=\E_{d_0}[s_0s_0^\top].
\]
因此，可观测矩为
\[
\widehat m_\phi=\frac1{2n_E}\sum_{i=1}^{n_E}z_i\bigl(t_\phi(e_i)-h_i\bigr).
\]
经验 ProRM 目标采用岭正则化。令
\[
\widehat\mu_F=\frac{1}{d}\operatorname{tr}(\widehat F_0),\qquad
\widehat F_{\lambda}=\widehat F_0+\lambda_{\rm rel}\widehat\mu_F I,
\]
则
\begin{equation}
\boxed{
\min_\phi\max_v\ \frac1\beta
\left[v^\top\widehat m_\phi-\frac12v^\top\widehat F_{\lambda}v\right].
}
\label{eq:minmax}
\end{equation}
在每个外层步骤中，先利用全部边计算 $\widehat m_\phi$，再以 PCG 求解
\[
\widehat F_{\lambda}v=\widehat m_\phi,
\]
随后将 $v$ 脱离计算图，依据包络定理给出的梯度更新奖励模型。对固定 $\phi$，内层是凹二次问题，因而
\[
v^*(\phi)=\widehat F_{\lambda}^{-1}\widehat m_\phi,\qquad
\mathcal L_{\mathrm{profile}}(\phi)=
\frac1{2\beta}\widehat m_\phi^\top\widehat F_{\lambda}^{-1}\widehat m_\phi.
\]
虽然 \cref{eq:minmax} 写成极小—极大形式，锁定实现并非同时执行小批量梯度下降/上升。每次外层更新均以全数据矩重新通过 PCG 求得近精确的 $v^*$，然后对 $\phi$ 执行一次 AdamW 下降。小批量仅用于等价的梯度累积与显存控制，不改变全数据梯度。直接对 $v$ 执行随机上升可作为近似鞍点算法，但会引入对偶学习率、内层滞后与收敛误差，故本实验不采用。当 reward head 为线性模型时，上述剖面目标是凸二次问题；BT-RM 则对应凸 logistic regression。

此处的 population 理论对象仍为 $F_0^\dagger$。$\widehat F_{\lambda}^{-1}$ 只是有限样本下的稳定化估计，不应解释为新的 population estimand。正则化强度必须仅依据训练数据选择，且选定后对所有随机种子和奖励来源共享；否则，方法特定的几何会与奖励质量发生混杂。

\begin{warnbox}[理论与有限样本实现的边界]
Population 理论使用 $F_0^\dagger$；工程实现使用 train-only 的
$\widehat F_\lambda^{-1}$。后者是稳定经验逆的估计方案，不是对 population
estimand 的重新定义。测试集只负责冻结后的评价，不进入 reward 拟合、Fisher
选择或方向求解。
\end{warnbox}

\subsection{验证问题与冻结设计}

本工程是 exact-$\Delta r^*$ 的受控机制实验：固定 prompt、六候选池、操作性
oracle、reward feature 与 rank-4 LoRA-$B$ tangent，只改变 MLE-RM 与 ProRM 的
训练目标。正式比较使用同一个 $\widehat F_\lambda$ 和同一个 $\beta$；oracle-NGD
是机制参照，tabular policy 是该有限候选池上的闭式最优上界。

\begin{table}[ht]
\centering
\small
\setlength{\tabcolsep}{4.2pt}
\begin{tabularx}{\textwidth}{>{\raggedright\arraybackslash\bfseries\color{navy}}p{3.35cm} X}
\toprule
\rowcolor{lightblue}
对象 & 冻结设置\\
\midrule
Seeds 与数据隔离 & Seeds 20261001--20261003；train/validation/test 为 3072/512/512。train 用于拟合、Fisher 与方向，test 仅用于本节评价。\\
Prompt 与候选 & MultiPref 只提供 prompts；$\pi_0=$ Qwen2.5-1.5B-Instruct；每个 prompt 有 $M=6$ 个独立生成候选。\\
Operational oracle & Skywork-Reward-V2-Llama-3.2-3B 的 scalar score 经 train-only 稳健仿射变换得到 $r^*$；它不是潜在人类效用。\\
Reward class & 冻结 Qwen response feature 加无 bias 的 1536 维线性 head。\\
Policy tangent & 最后一层 \texttt{q\_proj}/\texttt{v\_proj} 的 rank-4 LoRA-$B$，维度 7168；LoRA-$A$ 与 backbone 固定。\\
Fisher & Train nodes 的 raw second moment；五折 train-only cross-fit 从 $\lambda_{\rm rel}\in\{0.1,1,10\}$ 选择共享值 10。\\
主报告 $\beta$ & $\{0.1,0.2,0.3\}$；$0.2$ 为中心条件，$0.1/0.3$ 为两侧敏感性条件。\\
\bottomrule
\end{tabularx}
\caption{Main Experiment v1 的冻结对象。模型身份、数据顺序、精度和上游哈希均由 receipt 与 integrity audit 约束。}
\label{tab:engineering-config}
\end{table}

\subsection{Fisher 修正与共同正则系数下的单步 NGD}

令 $S\in\R^{N\times d}$ 的行是 train nodes 上的 score，定义
\begin{equation}
\widehat F=\frac1N S^\top S,\qquad
\widehat\mu_F=\frac1d\operatorname{tr}(\widehat F),\qquad
\widehat F_\lambda=\widehat F+\lambda_{\rm rel}\widehat\mu_F I.
\label{eq:empirical-fisher}
\end{equation}
这里的 damping 只服务于稳定求逆。对 reward source
$m\in\{\mathrm{MLE},\mathrm{Pro},\mathrm{oracle}\}$，先在 train 上冻结
\begin{equation}
\widehat g_m=A_{\rm train}\widehat r_m,\qquad
d_m=\widehat F_\lambda^{-1}\widehat g_m,\qquad
\delta_m(\beta)=\frac1\beta d_m.
\label{eq:ngd-step}
\end{equation}
因此该更新是固定正则系数下的一步 damped NGD，而不是把 KL 约束取等号后再校准的
TRPO。三种方向共享同一 Fisher 估计与阻尼，$\beta$ 只控制共同更新尺度。

在 test prompt $x_p$ 的六个冻结候选上，令 $s_{pi}$ 为候选 $i$ 的 policy score，
则四个参数化 policy 与 tabular optimum 分别为
\begin{align}
\pi_m(i\mid x_p)
&=\frac{\exp(s_{pi}^{\top}d_m/\beta)}{\sum_j\exp(s_{pj}^{\top}d_m/\beta)},
&&m\in\{\mathrm{MLE},\mathrm{Pro},\mathrm{oracle}\},\\
\pi_{\rm tab}(i\mid x_p)
&=\frac{\pi_0(i\mid x_p)\exp(r^*_{pi}/\beta)}
{\sum_j\pi_0(j\mid x_p)\exp(r^*_{pj}/\beta)}.
\label{eq:candidate-policies}
\end{align}
本实验候选由 $\pi_0$ 独立生成，有限池上的 reference 权重为 $1/M$。

\subsection{唯一正式评价指标}

Reward 层只保留三个诊断量。令 $q^*_{pe}=\sigma(\Delta r^*_{pe})$、
$q^m_{pe}=\sigma(\Delta\widehat r^m_{pe})$，则
\begin{align}
\operatorname{NLL}_m
&=-\frac1{|\mathcal E_{\rm test}|}\sum_{p,e}
\left[q^*_{pe}\log q^m_{pe}+(1-q^*_{pe})\log(1-q^m_{pe})\right],\\
\operatorname{MSE}_m
&=\frac1{PM}\sum_{p,i}
\left[(\widehat r^m_{pi}-\bar{\widehat r}^m_p)-(r^*_{pi}-\bar r^*_p)\right]^2,\\
\widetilde{\Reg}_{m,\beta}
&=\frac1{2\beta}\Delta A_{m,{\rm test}}^\top
\widehat F_{\lambda,{\rm train}}^{-1}\Delta A_{m,{\rm test}},
\quad
\Delta A_{m,{\rm test}}=A_{\widehat r_m,{\rm test}}-A_{r^*,{\rm test}}.
\label{eq:formal-reward-metrics}
\end{align}
NLL 与 MSE 不依赖 $\beta$；近似 regret 按 $1/\beta$ 精确缩放。它使用 test reward
moment 评价泛化，但保留训练阶段已冻结的 Fisher，不在 test 上重新选择几何。

Policy 层对每个 $\pi$ 定义
\begin{align}
R(\pi)&=\frac1P\sum_p\sum_i\pi(i\mid x_p)r^*_{pi},\\
K(\pi)&=\frac1P\sum_p\KL\!\left(\pi(\cdot\mid x_p)\,\|\,\pi_0(\cdot\mid x_p)\right),\\
J_\beta(\pi)&=R(\pi)-\beta K(\pi),\\
\Delta J_\beta(\pi)&=J_\beta(\pi_{\rm tab})-J_\beta(\pi),\\
\beta K_{\rm tab}(\pi)&=\frac\beta P\sum_p
\KL\!\left(\pi(\cdot\mid x_p)\,\|\,\pi_{\rm tab}(\cdot\mid x_p)\right),\\
J_{\rm close}(\beta)&=\frac\beta P\sum_p\log\left[
\frac1M\sum_i\exp(r^*_{pi}/\beta)\right].
\label{eq:formal-policy-metrics}
\end{align}
有限候选池上有两个严格恒等式
\begin{equation}
\boxed{J_\beta(\pi_{\rm tab})=J_{\rm close}(\beta)},\qquad
\boxed{\Delta J_\beta(\pi)=\beta K_{\rm tab}(\pi)}.
\label{eq:gibbs-identities}
\end{equation}
因此 $\Delta J$ 与 $\beta K_{\rm tab}$ 是同一 gap 的两种独立实现，用于交叉校验，
不是两项独立证据。

\subsection{结构化工程链路与可接续边界}

\begin{figure}[ht]
\centering
\begin{tikzpicture}[
  node distance=7mm and 7mm,
  box/.style={rounded corners=1.2mm,draw=line,fill=paperblue,align=center,
    minimum height=13mm,text width=3.2cm,font=\sffamily\footnotesize},
  arr/.style={-Latex,thick,draw=navy}
]
\node[box] (anc) {{\color{navy}\bfseries 01 验证上游遗产}\\material · reward · Fisher};
\node[box,right=of anc] (eval) {{\color{navy}\bfseries 02 三 seed evaluation}\\五 policies · dense $\beta$};
\node[box,right=of eval] (agg) {{\color{navy}\bfseries 03 Aggregate}\\seed 为统计单位};
\node[box,below=of agg] (audit) {{\color{navy}\bfseries 04 Integrity audit}\\身份 · 哈希 · 恒等式};
\node[box,left=of audit] (archive) {{\color{navy}\bfseries 05 Immutable archive}\\HPC4 项目盘};
\node[box,left=of archive] (report) {{\color{navy}\bfseries 06 Main v1 bundle}\\Git · 本地数据 · PDF};
\draw[arr] (anc)--(eval);
\draw[arr] (eval)--(agg);
\draw[arr] (agg)--(audit);
\draw[arr] (audit)--(archive);
\draw[arr] (archive)--(report);
\end{tikzpicture}
\caption{本轮主实验的依赖闭包。上游 artifact、reward heads、Fisher 与方向均未重算；只扩展共同 $\beta$ evaluation 及其下游聚合、审计与报告。}
\label{fig:pipeline}
\end{figure}

每个阶段的 immutable receipt 与内容哈希决定能否接续。只有最早受影响阶段及其科学下游
需要重算；没有受影响且验证通过的上游不重复生成。最终 Git bundle 同时保留三个 seed
的完整 evaluation JSON、aggregate、audit、archive manifest、机器可读 summary 与
由 summary 确定性生成的结果表。

\subsection{Main Experiment v1 结果}

\InputIfFileExists{ProRM_main_results.tex}{}{%
  \begin{warnbox}[结果文件缺失]
  请先运行 \texttt{scripts/reporting/build\_ngd\_main\_report.py}，从通过审计的
  aggregate 与三个 seed evaluation 生成本节表格。
  \end{warnbox}
}

\begin{claimbox}[下游 policy 结论]
在 $\beta=0.1,0.2,0.3$ 的三 seed 聚合均值上，均观察到
\[
R_{\rm tab}>R_{\rm oracle}>R_{\rm Pro}>R_{\rm MLE}>R_{\pi_0},
\qquad
J_{\rm tab}>J_{\rm Pro}>J_{\rm oracle}>J_{\rm MLE}>J_{\pi_0}.
\]
逐 seed 检查中，Pro 的 $J$ 在全部 $9/9$ 个 $\text{seed}\times\beta$ 条件上同时高于
MLE 与 oracle；完整 $J$ 链在 $\beta=0.1$ 为 $2/3$，在 $\beta=0.2,0.3$ 均为 $3/3$，
唯一偏离来自一个 seed 的 oracle 与 MLE 次序互换，不影响 Pro--MLE 主比较。
中心条件 $\beta=0.2$ 下，Pro 相对 MLE 的平均 $J$ 优势为 $0.005098$；Pro 与
oracle 的平均 $J$ 差为 $0.000935$。因此 Pro-NGD 明显优于 MLE-NGD，并保持接近
oracle-NGD 的 regularized utility。
\end{claimbox}

\begin{warnbox}[Reward 层负结果：RM 过拟合/泛化不足]
ProRM 的 train profiled objective 确实被优化：三个 seed 上均低于 MLE 投影，均值由
$0.004951$ 降至 $0.002115$。但在 held-out test 上，ProRM 的 NLL、centered reward
MSE 与近似 regret 全部差于 MLE-RM；两个 seed 的 Pro head 范数还从 MLE 的约
$4.94/3.55$ 增长到 $11.28/11.69$。这说明当前主要问题是 reward head 对训练矩目标
过拟合、policy-relevant moment 泛化不足，而不是 PCG、Fisher 求逆或 NGD 公式失效。
该问题必须诚实保留为限制；本轮不以事后加正则化改写结果。
\end{warnbox}

\begin{warnbox}[证据等级与外部有效性]
$\beta=0.1,0.2,0.3$ 是观察 dense test sweep 后确定的主报告区间，因此本版属于
完整、可审计的探索性主实验，而不是新 held-out confirmatory test。未来确认性实验应在
新 test 前冻结 $\beta=0.2$（$0.1/0.3$ 为敏感性）。此外，policy 指标来自冻结六候选池，
不直接等价于多步 PPO/GRPO 或开放式 fresh generation；Skywork 也只是 operational
oracle。三个 seed 支持一致性与描述统计，不支持渐近显著性主张。
\end{warnbox}

\Needspace{7\baselineskip}
\begin{flushleft}\footnotesize\color{gray}
Producer commit: \texttt{0f51c6a59454e7ad7e4fa9fd09f022feb0598d68}\\
Aggregate SHA-256: \texttt{fb6967abf4a8b0ca679c6a022c9537a6d1dc471a90187580e110b751605ac930}\\
Audit status: \texttt{passed}; archive manifest SHA-256:
\texttt{fac001daf941a365d6db4541219e24b26d01e028d9a158e111a092040d097e85}.
\end{flushleft}

% ================================================================
\section*{讨论与结论}
\addcontentsline{toc}{section}{讨论与结论}

ProRM 的理论主线是：奖励模型不应只解释偏好概率，而应最小化它所诱导的下游 policy regret。Taylor--Fisher 局部化把双层效用转化为 policy-visible quotient space 中的投影范数；奖励等价类进一步说明，完整恢复 $r^*$ 并非必要。Main Experiment v1 修正了经验 Fisher 的估计与正则逆，并用共同 $\beta$ 的单步 NGD 对五个 policy 进行有限候选池评价。结果呈现出一个必须同时报告的张力：ProRM 的 reward-level test 泛化弱于 MLE-RM，但 Pro-NGD 的下游 $J$ 在三个 $\beta$、三个 seed 上均优于 MLE-NGD，并接近 oracle-NGD。因而本轮最强结论不是“ProRM 在所有 reward 指标上更好”，而是“即使传统 reward-fit 与局部 proxy 更差，面向决策的训练目标仍可产生更优的下游 regularized utility”。这正是本文理论所强调的评价错位，同时也指出下一轮应直接处理 reward-moment 泛化，而不应掩盖该限制。

\enlargethispage{0.8\baselineskip}
\begin{multicols}{2}
\setlength{\columnsep}{4mm}
\begin{thebibliography}{99}
\fontsize{6.8}{6.9}\selectfont
\setlength{\itemsep}{-0.25em}
\setlength{\parsep}{0pt}
\bibitem{spo}
A. N. Elmachtoub and P. Grigas. Smart ``Predict, then Optimize''. arXiv:1710.08005, 2020 revision.
\bibitem{bt}
R. A. Bradley and M. E. Terry. Rank analysis of incomplete block designs: I. The method of paired comparisons. \emph{Biometrika}, 39(3/4):324--345, 1952.
\bibitem{amari}
S. Amari. \emph{Information Geometry and Its Applications}. Springer, 2016.
\bibitem{trpo}
J. Schulman, S. Levine, P. Abbeel, M. Jordan, and P. Moritz. Trust Region Policy Optimization. \emph{ICML}, PMLR 37, 2015.
\bibitem{zhu_principled}
B. Zhu, M. I. Jordan, and J. Jiao. Principled Reinforcement Learning with Human Feedback from Pairwise or K-wise Comparisons. arXiv:2301.11270, 2024.
\bibitem{wen_rm_eval}
X. Wen, J. Lou, Y. Lu, et al. Rethinking Reward Model Evaluation: Are We Barking Up the Wrong Tree? \emph{ICLR}, 2025.
\bibitem{nika_rml_dpo}
A. Nika, D. Mandal, P. Kamalaruban, G. Tzannetos, G. Radanovi\'c, and A. Singla. Reward Model Learning vs.\ Direct Policy Optimization: A Comparative Analysis of Learning from Human Preferences. \emph{ICML}, PMLR 235, 2024.
\bibitem{li_rl_dpo}
Z. Li, T. Xu, and Y. Yu. When Is RL Better Than DPO in RLHF? A Representation and Optimization Perspective. \emph{ICLR Tiny Papers}, 2024.
\bibitem{shi_gap}
R. Shi, M. Song, R. Zhou, Z. Zhang, M. Fazel, and S. S. Du. Understanding the Performance Gap in Preference Learning: A Dichotomy of RLHF and DPO. \emph{ICML}, PMLR 306, 2026.
\bibitem{yang_equivalence}
Z. Yang, Y. Zhang, W. Xue, D. Fang, B. Han, and Y. Guo. Conditional Equivalence of DPO and RLHF: Implicit Assumption, Failure Modes, and Provable Alignment. \emph{ICML}, PMLR 306, 2026.
\bibitem{han_pairs}
J. Han, V. Goyal, and W. Ma. Which Pairs to Compare for LLM Post-Training? arXiv:2606.19607, 2026.
\bibitem{gopalan_misspecified}
A. Gopalan, S. Ray Chowdhury, and D. Banerjee. Why DPO Is a Misspecified Estimator and How to Fix It. arXiv:2510.20413, 2025.
\end{thebibliography}
\end{multicols}

\end{document}
